\section{Reduction to fixed-precision approximation}
\label{sec:deduction}

The main theorems follow from an approximation algorithm with a fixed
inverse-polynomial error bound. We state that intermediate result here
and give the deductions. Its construction, including the median and
the simultaneous time and space bounds, is proved in
Section~\ref{sec:conclusion}.

All transducers have read-only input and a one-way write-only output
tape. For a fixed randomized machine $\mathcal M$, write
$p_{\mathcal M}(x)$ for its acceptance probability on $x$.

\begin{theorem}[Fixed inverse-polynomial accuracy]
\label{thm:fixed-precision}
Fix a randomized machine $\mathcal M$ with polynomial worst-case
running time and $O(\log(n+2))$ work space, and fix an integer $d\ge1$.
There is a uniform deterministic transducer that, on every input $x$
of length $n$, outputs a dyadic rational $z_x\in[0,1]$ with
\[
 |z_x-p_{\mathcal M}(x)|\le(n+2)^{-d}.
\]
Its work space and output length are $O(\log(n+2))$, and its running
time is polynomial in $n+2$. The constants may depend on
$\mathcal M$ and $d$, but not on the input.
\end{theorem}

The proof in Section~\ref{sec:conclusion} combines the deterministic
reward approximation of Proposition~\ref{prop:hier-reward-target},
the moment bound of Theorem~\ref{thm:fp-estimators}, and the total
space-bounded implementation of Theorem~\ref{thm:ctrl-realization}.
The deductions below need only the stated guarantee of
Theorem~\ref{thm:fixed-precision}.

\subsection{The class equality and promise separation}

\begin{proof}[Proof of Theorem~\ref{thm:main}]
Let $\mathcal M$ recognize a language in $\BPL$, with acceptance
probability at most $1/3$ on no inputs and at least $2/3$ on yes
inputs. Apply Theorem~\ref{thm:fixed-precision} with $d=3$.
For every input length, including $n=0$, its error is at most
$(n+2)^{-3}\le1/8$. Hence the dyadic answer is at most
$1/3+1/8=11/24$ on a no input and at least
$2/3-1/8=13/24$ on a yes input. Comparing this answer with $1/2$
decides the language uniformly in logarithmic space. This proves
$\BPL\subseteq\Lclass$ directly from probability approximation.

A deterministic machine may ignore its coins, giving
$\Lclass\subseteq\RL$. Two independent repetitions of an $\RL$
machine, accepting if either repetition accepts, have acceptance
probability zero on no inputs and at least $3/4$ on yes inputs.
They use logarithmic space and polynomial time, so
$\RL\subseteq\BPL$. The three inclusions prove
$\Lclass=\RL=\BPL$.
\end{proof}

A promise problem is a pair $(Y,N)$ of disjoint sets of inputs.
The class $\mathsf{prBPL}$ of the introduction consists of the pairs for
which one polynomial-time randomized
logarithmic-space machine halts on every input, accepts each
input in $Y$ with probability at least $2/3$, and accepts each
input in $N$ with probability at most $1/3$.
The class $\mathsf{prL}$ consists of pairs separated by a total
deterministic logarithmic-space decider; there is no correctness
requirement outside $Y\cup N$.

\begin{corollary}\label{cor:promise-bpl}
Under these conventions,
$\mathsf{prBPL}=\mathsf{prL}$.
\end{corollary}

\begin{proof}
The same $d=3$ approximation and threshold $1/2$ used in the preceding
proof accept every input in $Y$ and reject every input in $N$.
The resulting machine halts also outside the promise because its
approximation queries are total. The reverse inclusion follows by
ignoring the coins.
\end{proof}

\subsection{Accuracy supplied as part of the input}

Theorem~\ref{thm:fixed-precision} fixes the accuracy exponent before
reading the input.  To obtain one algorithm for arbitrary input precision
$q$, we instead fix its exponent to one and apply it to a longer input
represented implicitly.  Padding preserves the source acceptance
probability while making the inverse-linear error small enough for the
requested accuracy.

\begin{proof}[Proof of Theorem~\ref{thm:probability}]
Fix a padded machine $\mathcal M_{\rm pad}$ as follows. On a string
beginning with $1^n0x$, where $|x|=n$, it simulates $\mathcal M(x)$
and ignores all remaining symbols. If the delimiter or the following
$n$ symbols are absent, it rejects. This is one fixed randomized
machine with polynomial worst-case running time and logarithmic work
space in its full input length. On every string with the specified
prefix its acceptance probability is $p_{\mathcal M}(x)$.
Fix once and for all the transducer of
Theorem~\ref{thm:fixed-precision} for $\mathcal M_{\rm pad}$ and $d=1$.

On input $(x,1^q)$, with $n=|x|$ and $q\ge1$, set
\[
 m=2^{q+\lceil\log_2(2n+4)\rceil+4}.
\]
Simulate that fixed transducer on the virtual length-$m$ string
\[
 y=1^n0x0^{m-(2n+1)}.
\]
An input query is answered by comparing its address with the prefix
boundaries and either returning the appropriate bit of $x$ or a
fixed bit. The simulated input-head address and all these comparisons
use $O(\log m)$ bits; the string $y$ is never stored. The simulated
transducer returns a dyadic rational $z\in[0,1]$ satisfying
\[
 |z-p_{\mathcal M}(x)|\le(m+2)^{-1}<2^{-q-4}.
\]
Its numerator and exponent have $O(\log m)$ bits and can be retained
in work space while its finite output is simulated.

Compute by exact dyadic arithmetic
\[
 a=\left\lfloor 2^{q+2}z+\tfrac12\right\rfloor.
\]
Then $0\le a\le2^{q+2}$, and rounding gives
\[
 \left|a2^{-(q+2)}-p_{\mathcal M}(x)\right|
 \le2^{-q-3}+2^{-q-4}<2^{-q}.
\]
Write the binary numerator $a$ to the output tape. It has at most
$q+3$ bits and represents the denominator specified in the theorem.
The simulator state, virtual input address, actual input counters,
captured dyadic answer, and rounding operands occupy a fixed number
of records of total size
\[
 O(\log m)=O(\log(n+2)+q).
\]
The simulation and the synthesis of each input bit take polynomial
time in $m$, so the total time is
$m^{O(1)}=(n+2)^{O(1)}2^{O(q)}$.
The padded machine and its exponent-one approximator were fixed
before reading $q$. Thus their constants are independent of the
requested accuracy. Throughout their simulation, the original size
convention is $B=\Theta(\log(m+2))$, exactly as required by the
estimator and controller theorems.
\end{proof}

